\section{Methodology}

The computational experiments were performed using Python and the NumPy library.
All problem-specific data were generated deterministically from the student's
enrollment number, following the algorithms provided in the assignment.

\subsection{Matrix Generation and Cholesky Factorization}

For the first problem, the matrix $A\in\mathbb{R}^{10\times10}$ is generated according to

\begin{equation}
    A = BB^{T} + 10I,
\end{equation}

where $B$ is a randomly generated integer matrix and $I$ is the $10\times10$
identity matrix. The vector $b\in\mathbb{R}^{10}$ is independently generated
using the prescribed random seed.

The symmetry and positive definiteness of $A$ are verified before computing its
Cholesky factorization. When $A$ is symmetric positive definite, the
factorization

\begin{equation}
    A = GG^{T}
\end{equation}

is computed, where $G$ is a lower triangular matrix with positive diagonal
entries.

The condition number of $A$ is then computed to quantify the sensitivity of the
corresponding linear system to perturbations \cite{golub2013matrix}. Finally, the system

\begin{equation}
    Ax=b
\end{equation}

is solved numerically.

\subsection{Heat Conduction Problem}

For the second problem, the parameters $k$, $f$, and the boundary-condition type
are generated using the prescribed procedure and the enrollment number as the
random seed. The spatial domain $[0,1]$ is represented by $11$ equally spaced
grid points, resulting in a grid spacing

\begin{equation}
    h=\frac{1}{10}.
\end{equation}

The matrix $M$ and vector $s$ are constructed according to the finite difference
formulation specified in the assignment. The resulting linear system

\begin{equation}
    MT=s
\end{equation}

is then solved numerically to obtain the temperature vector

\begin{equation}
    T =
    \begin{bmatrix}
        T_0 & T_1 & \cdots & T_{10}
    \end{bmatrix}^{T}.
\end{equation}

The numerical temperature values are associated with the corresponding spatial
grid points and plotted as a function of $x$. The resulting temperature
distribution is subsequently interpreted in terms of the physical behavior of
the one-dimensional heat conduction problem.

All numerical computations were performed using double-precision floating-point
arithmetic. For convergence analysis, the discrete $L_2$ error is explicitly defined as
\begin{equation}
    \|e\|_{2,h} = \sqrt{h} \, \|e\|_2,
\end{equation}
which is the exact formulation implemented in the computational experiment.